(1) 令 \(x=h\nu/(k_BT)\)，有 \(d\nu=(k_BT/h)dx\)。积分得 \(u(T)=8\pi(k_BT)^4/(h^3c^3)\int_0^\infty x^3/(e^x-1)dx=8\pi^5 k_B^4T^4/(15h^3c^3)\)。因此光子能量密度严格按 \(T^4\) 增长，而不依赖容器形状。
(2) 在各向同性辐射场中，取半球角积分 \(\int_{\rm hemi}\cos\theta\,d\Omega/(4\pi)=1/4\)，向外能流 \(F=cu/4=\sigma T^4\)，其中 \(\sigma=2\pi^5k_B^4/(15h^3c^2)\)。利用 \([u]=\mathrm{J\,m^{-3}}\) 可验 \([F]=\mathrm{W\,m^{-2}}\)。
(3) 谱变换须用能量守恒 \(B_\lambda d\lambda=B_\nu|d\nu|\)，且 \(|d\nu/d\lambda|=c/\lambda^2\)；于是 \(B_\lambda\propto\lambda^{-5}/(e^{hc/(\lambda k_BT)}-1)\)。令 \(x=hc/(\lambda k_BT)\)，峰值条件为 \(5(1-e^{-x})=x\)，正根 \(x\simeq4.9651\)，从而 \(\lambda_{\max}T=hc/(4.9651\,k_B)\)。
(4) 频率谱 \(B_\nu\) 的峰值根满足 \(3(1-e^{-x_\nu})=x_\nu\)，与波长谱的系数 \(5\) 不同。因为变量替换时密度还乘以 Jacobian，单点谱强度并不保持不变；故 \(c/\nu_{\max}\) 不等于 \(\lambda_{\max}\)，但对全谱积分后的总辐射功率一致。
